\documentclass{article}

\usepackage[a4paper,landscape,margin=2cm]{geometry}
\usepackage{../../../../components/mathtex}
\usepackage{colortbl}
\usepackage{array}

\usepackage{fancyhdr}
\usepackage{ulem}
\pagestyle{fancy}
\fancyhf{}
\fancyhead[L]{DL3 Math-Info}
\fancyhead[C]{Fiche d'entraînement hebdomadaire N°4}
\fancyhead[R]{2026-2027}
\fancyfoot[L]{Ewen Rodrigues de Oliveira}
\fancyfoot[R]{\thepage}

\usepackage{paracol}

\begin{document}
\setlength{\columnsep}{2cm}

\begin{paracol}{2}

\renewcommand{\arraystretch}{1.3}
\noindent
\begin{tabular}{|p{0.6\columnwidth}|>{\centering\arraybackslash}p{0.35\columnwidth}|}
\hline
\rowcolor[RGB]{230,237,247}
\textbf{Matière} & \textbf{Exercices} \\
\hline
Intégration et probabilités & $10,11,12,13,14,15$\\
\hline
Groupes et actions & $1,2,3,4,5,6$\\
\hline
Calcul différentiel & $7,8,9,16,17$\\
\hline
Systèmes d'exploitation & $\emptyset$\\
\hline
Algorithmique & $\emptyset$\\
\hline
Programmation fonctionnelle & $\emptyset$\\
\hline
Anglais & $\emptyset$\\
\hline
\end{tabular}

\vspace{1em}
\exercise{1}{Propriétés des groupes finis}{
    Soit $(G,*)$ un groupe fini.
    \begin{enumerate}
        \item On suppose que $G$ est d'ordre pair. Démontrer que le cardinal de l'ensemble $\{g \in G \mid g = g^{-1}\}$ est pair et non nul.
        \item On suppose que l'application $G \to G, \; x \mapsto x^2$ est bijective. Démontrer que $G$ est d'ordre impair.
    \end{enumerate}
}

\exercise{2}{Sous-groupes de $\mathbb{Z}/n\mathbb{Z}$}{
    Soit $n \in \mathbb{N}\setminus\{0\}$.
    \begin{enumerate}
        \item Démontrer que les groupes $\mathbb{Z}/n\mathbb{Z}$ et $\mathbb{U}_n$ sont isomorphes.
        \item En déduire que les seuls sous-groupes de $\mathbb{Z}/n\mathbb{Z}$ sont : $k\mathbb{Z}/n\mathbb{Z}$ ($k$ divisant $n$).
    \end{enumerate}
}

\exercise{3}{Homomorphismes entre groupes d'ordres premiers entre eux}{
    Soient $G$ et $G'$ deux groupes finis. On suppose que les ordres de ces groupes sont premiers entre eux. Déterminer tous les homomorphismes de $G$ dans $G'$.
}

\switchcolumn

\exercise{4}{Image d'un sous-groupe par un morphisme}{
    Soit $f : G \to H$ un morphisme de groupes finis et $K$ un sous-groupe de $G$.
    \begin{enumerate}
        \item Montrer que l'ordre de $f(K)$ divise les ordres de $K$ et de $H$.
        \item Montrer que si l'ordre de $K$ est premier à l'ordre de $H$ alors $K \subset \mathrm{Ker}(f)$.
    \end{enumerate}
}

\exercise{5}{Isomorphismes et morphismes de groupes}{
    \begin{enumerate}
        \item Le groupe $(\mathbb{R}^*, \times)$ est-il isomorphe au groupe $(\mathbb{C}^*, \times)$ ?
        \item Le groupe $(\mathbb{C}, +)$ est-il isomorphe au groupe $(\mathbb{C}^*, \times)$ ?
        \item Trouver tous les endomorphismes de $(\mathbb{Q}, +)$.
        \item Trouver tous les morphismes de $\mathbb{Z}$ dans $\mathbb{Q}$.
    \end{enumerate}
}

\exercise{6}{Étude d'un groupe de matrices}{
    Soit $G$ l'ensemble des matrices de la forme :
    \[
    \begin{pmatrix} 1 & 0 \\ a & \varepsilon \end{pmatrix}
    \]
    avec $a \in \mathbb{Z}$ et $\varepsilon \in \{-1, 1\}$. On pose $r = \begin{pmatrix} 1 & 0 \\ 1 & 1 \end{pmatrix}$ et $s = \begin{pmatrix} 1 & 0 \\ 0 & -1 \end{pmatrix}$.
    \begin{enumerate}
        \item 
        \begin{enumerate}
            \item Montrer que $G$ est un sous-groupe de $(\mathrm{GL}_2(\mathbb{R}), \times)$.
            \item Est-il abélien ?
        \end{enumerate}
        \item 
        \begin{enumerate}
            \item Montrer que l'application
            \[
            f : \mathbb{Z} \to G, \quad k \mapsto \begin{pmatrix} 1 & 0 \\ k & 1 \end{pmatrix}
            \]
            est un morphisme de groupes de $(\mathbb{Z}, +)$ dans $G$.
            \switchcolumn
            \newpage
            \item Vérifier que $\mathrm{Im}(f) = \langle r \rangle$.
            \item Montrer que $\langle r \rangle$ est isomorphe à $\mathbb{Z}$ puis que $r^k$ est d'ordre infini pour chaque entier relatif non nul $k$.
            \item Montrer que, pour tout sous-groupe $H$ de $G$, il existe un entier naturel $n$ tel que $H \cap \langle r \rangle = \langle r^n \rangle$.
        \end{enumerate}
        \item 
        \begin{enumerate}
            \item Calculer $sr^k$ pour $k \in \mathbb{Z}$ et justifier que $G$ est engendré par $r$ et $s$.
            \item Quel est l'ordre de $sr^k$ ?
        \end{enumerate}
        \item Soient $m$ et $n$ deux entiers relatifs.
        \begin{enumerate}
            \item Montrer qu'il existe un unique morphisme de groupes $\varphi$ de $G$ dans $G$ tel que $\varphi(r) = r^m$ et $\varphi(s) = sr^n$. On pourra commencer par supposer l'existence d'un tel morphisme et exprimer pour chaque entier $k$ les éléments $\varphi(r^k)$ et $\varphi(sr^k)$ en fonction de $k$.
            \item Préciser pour quelles valeurs du couple $(m, n)$ on obtient un automorphisme de $G$.
        \end{enumerate}
    \end{enumerate}
}

\exercise{7}{Normes}{
    Soient $N_1, N_2 : \mathbb{R}^n \to \mathbb{R}$ deux normes. On définit $N : \mathbb{R}^n \to \mathbb{R}$ par $N(x) = N_1(x) + 2N_2(x)$. Montrer que $N$ est une norme.
}

\exercise{8}{Topologie}{
    Pour chacun des ensembles suivants, déterminer leur intérieur et leur adhérence. Sont-ils ouverts ? Fermés ? \textit{On n'attend pas de justification pour cet exercice.}
    \begin{itemize}
        \item $A = [-1, 1[$ un ensemble de $\mathbb{R}$.
        \item $B = \overline{\mathbb{B}}(0, 2) \setminus \mathbb{B}(0, 1)$ un ensemble de $\mathbb{R}^2$. \textit{Faire un dessin vous aidera certainement.}
    \end{itemize}
}

\exercise{9}{Topologie}{
    Soit $I = [1, +\infty[ \subset \mathbb{R}$. Montrer que l'intérieur de $I$ est égal à $]1, +\infty[$. \textit{Une preuve soignée est attendue. Tout raisonnement bien rédigé bien qu'incomplet sera récompensé.}
}

\switchcolumn
\newpage

\exercise{10}{Question de cours}{
    Montrer que $\mathcal{B}(\mathbb{R}^d)$ est généré par $\{]-inf, a_i[ | a_i \in \mathbb{Q}\}$, où $\mathcal{B}(\mathbb{R}^d)$ est la tribu borélienne de $\mathbb{R}^d$.
}

\exercise{11}{Préparation du CM}{Passer le QCM "Rappels de cours 1" sur la page Moodle de Intégration et Probabilités.}
\exercise{12}{Questions de cours}{
    \textit{(on pourra, dans cet exercice uniquement, donner les réponses sans justification)}
    \begin{enumerate}
        \item Donner la définition de tribu sur un univers $\Omega$.
        \item Soient $k \in \mathbb{N}^*$, $\ell \in \mathbb{N}^*$, quel est le cardinal de $\{B \in \mathcal{P}(\{1,\dots,\ell\}) : \mathrm{card}(B) = k\}$ ?
        \item Soit $(\Omega, \mathcal{F}, \mathbb{P})$ un espace probabilisé. Pour un $n \in \mathbb{N}^*$ on se donne $\{A_1, A_2, \dots, A_n\}$ des éléments de $\mathcal{F}$ formant une partition de $\Omega$ et on suppose que $\mathbb{P}(A_i) > 0$, pour tout $i \in \{1,\dots,n\}$. Pour $B \in \mathcal{F}$, que vaut
        \[
        \sum_{k=1}^n \mathbb{P}(B \mid A_i)\mathbb{P}(A_i) \quad ?
        \]
    \end{enumerate}
}

\exercise{13}{Étude d'une mesure de probabilité sur $\mathbb{N}$}{
    On considère $\Omega = \mathbb{N}$ muni de la tribu de ses parties, et d'une mesure de probabilité $\mathbb{P}$ telle que
    \[
    \mathbb{P}(\{0\}) = \mathbb{P}(\{1\}), \quad \mathbb{P}(\{2\}) = 0, \quad \text{et}
    \]
    \[
    \mathbb{P}(\{k\}) = 3^{2-k}, \quad \forall k \in \mathbb{N}, k \ge 3.
    \]
    \begin{enumerate}
        \item Établir que
        \[
        \mathbb{P}(\mathbb{N} \setminus \{0, 1, 2\}) = \frac{1}{2}
        \]
        \item En déduire $\mathbb{P}(\{0\}), \mathbb{P}(\{1\})$.
        \item Pour tout $x \in (-3, 3)$ montrer que
        \[
        \sum_{k \in \mathbb{N}} x^k 3^{-k} = \frac{3}{3 - x}.
        \]
        \switchcolumn
        \newpage
        En déduire que
        \[
        \sum_{k \in \mathbb{N}, k \ge 3} k 3^{2-k} = \frac{7}{4},
        \]
        et calculer finalement
        \[
        \sum_{k \in \mathbb{N}} k \mathbb{P}(\{k\}).
        \]
        \item Calculer $\sum_{k \in \mathbb{N}} k^2 \mathbb{P}(\{k\})$.
    \end{enumerate}
}

\exercise{14}{Mesure de probabilité paramétrique}{
    Soit $\lambda > 0$. Existe-t-il une constante $c(\lambda)$ telle que
    \[
    \mathbb{P}(\{i\}) = c(\lambda) \frac{(i + 1)\lambda^i}{i!} \quad \text{pour tout } i \in \mathbb{N},
    \]
    définisse une probabilité sur $\mathbb{N}$ ?
}

\exercise{15}{Mesures définies par une densité}{
    Soit $f : \mathbb{R} \to \mathbb{R}_+$ une fonction positive, continue par morceaux. On admet qu'il existe une unique mesure positive $\mu$ sur $\mathcal{B}(\mathbb{R})$ telle pour tous $a < b$,
    \[
    \mu([a,b]) = \int_a^b f(x) dx.
    \]
    \begin{enumerate}
        \item Soit $a \in \mathbb{R}$, que vaut $\mu(\{a\})$ ?
        \item Peut-il y avoir une différence entre $\mu((a,b))$, $\mu((a,b])$ et $\mu([a,b])$ ?
        \item Que vaut $\mu(\mathbb{Q})$ ?
        \item Exprimer $\mu(\mathbb{R} \setminus ([-2, -1] \cup [1, 2]))$ en fonction de $f$. Calculer cette quantité lorsque $f(x) = \exp(-|x|), x \in \mathbb{R}$
        \item Exprimer $\mu\left(\bigcup_{n \in \mathbb{N}} [0, n]\right)$ en fonction de $f$. Calculer cette quantité lorsque $f(x) = \frac{1}{1+x^2}, x \in \mathbb{R}$.
        \item À quelle condition sur $f$ la mesure $\mu$ est-elle finie ? À quelle condition est-elle une probabilité ? Donner un exemple.
        \item Comment appelle-t-on $\mu$ lorsque $f \equiv 1$ ?
    \end{enumerate}
}

\switchcolumn
\exercise{16}{Topologie en dimension infinie}{
    On considère $E = \ell^\infty$ l'espace vectoriel des suites réelles bornées. Pour tout $u = (u(k))_k \in E$ on pose :
    \[
    \|u\|_\infty = \sup_{k \in \mathbb{N}} |u(k)|
    \]
    On considère deux sous-espaces vectoriels de $E$, définis par :
    \[
    A := \left\{ u = (u(k))_k \in E \;\middle|\; u(k) \xrightarrow[k \to +\infty]{} 0 \right\}
    \]
    \[
    B := \{ u = (u(k))_k \in E \mid \exists K \in \mathbb{N}, \forall k \ge K, u(k) = 0 \}
    \]
    \begin{enumerate}
        \item Montrer que $\|\cdot\|_\infty$ est une norme sur $E$.
        \item Montrer que $A$ est fermé dans $(E, \|\cdot\|_\infty)$, mais que $B$ ne l'est pas.\\
        \textit{Indication : On pourra utiliser la suite $(u_n)_n$ d'éléments de $B$ définie par $u_n(k) = \begin{cases} 1 & \text{si } k \le n \\ 0 & \text{sinon} \end{cases}$.}
        \item Montrer que $A$ est l'adhérence de $B$ dans $(E, \|\cdot\|_\infty)$.
        \item Parmi les espaces $(E, \|\cdot\|_\infty)$, $(A, \|\cdot\|_\infty)$ et $(B, \|\cdot\|_\infty)$, lesquels sont complets ?
    \end{enumerate}
}

\exercise{17}{Étude de l'espace $\ell^1$}{
    Soit $\ell^1$ l'espace vectoriel des applications $u : \mathbb{N} \to \mathbb{R}$ pour lesquelles la somme $\sum_n |u(n)|$ converge. Pour chaque élément $u$ de $\ell^1$, on note
    \[
    \|u\|_1 = \sum_n |u(n)|.
    \]
    \begin{enumerate}
        \item Démontrer que $\|\cdot\|_1$ est une norme sur $\ell^1$.\\
        \textit{Dans la suite de cette exercice, on munit l'espace vectoriel $\ell^1$ de la norme $\|\cdot\|_1$.}
        \switchcolumn
        \item Soit $u$ un élément de $\ell^1$, et considérons les éléments $u_k$ de $\ell^1$ tels que $u_k(n) = u(n)$ pour $n \le k$ et $u_k(n) = 0$ pour $n > k$. Démontrer que la suite $(u_k)_k$ converge vers $u$ dans l'espace vectoriel normé $\ell^1$.
        \item Démontrer que l'ensemble $F$ des éléments $u$ de $\ell^1$ tels que $\sum_n u(n) = 0$ est une partie fermée de $\ell^1$.
        \item Déterminer l'intérieur de $F$.
    \end{enumerate}
}

\end{paracol}

\end{document}
